\clearpage\begingroup\SourceLeaf{366b}\ArithmeticHead{35}{R}{DIE GRUNDLAGEN DER GEOMETRIE.}
\renewcommand{\thefootnote}{\arabic{footnote})}\setcounter{footnote}{0}
\begin{center}
\textbf{8.}\par\smallskip
\textls[100]{Die weitere Entwicklung bei \textsc{Gauss};}\\
\textls[100]{\textsc{Johann Bolyai} und \textsc{Lobatschefskij} (1831—1846).}
\end{center}
\noindent\hspace*{\parindent}Am 3. November 1823 hatte \textsc{Johann Bolyai} aus Temesvar, wo er als
Pionierleutnant stand, seinem Vater mitgeteilt, er habe »aus Nichts eine neue,
andere Welt geschaffen«. Im Februar 1825 konnte er ihm den ersten Entwurf
seiner absoluten Raumlehre vorlegen. \textsc{Wolfgang} war jedoch damit nicht
einverstanden; besonders nahm er Anstoss an dem Auftreten der an sich unbestimmten
Konstanten und der dadurch bedingten Vielheit der möglichen
hypothetischen Systeme (\textsc{Bol.} S. 87). Vater und Sohn konnten sich nicht
\par\vfill\hfill{\footnotesize 5*}\hspace*{2em}
\clearpage\SourceLeaf{367}\ArithmeticHead{36}{L}{STÄCKEL, GAUSS ALS GEOMETER.}\setcounter{footnote}{0}
\noindent einigen, und schliesslich kam man überein, \textsc{Johann} möge das Wesen der Sache
in lateinischer Sprache darstellen, die kleine Abhandlung solle dem von \textsc{Wolfgang}
geplanten \textit{Tentamen}\footnote{W. \textsc{Bolyai}, \textit{Tentamen iuventutem studiosam in elementa matheseos .\ .\ .\ introducendi}, t. I, Maros Vásárhely 1832, ed. secunda, Budapest 1897.} beigegeben und einer der herzustellenden Abzüge
an \textsc{Gauss} gesandt werden; seinem Urteil über Wert oder Unwert wollten sich
beide unterwerfen.
\par
Im Juni 1831 wurden die Sonderabzüge des \textit{Appendix scientiam spatii absolute
veram exhibens} fertig, und am 20. Juni wurde einer davon an \textsc{Gauss}
abgesandt. Jedoch gelangte »der fatalen Choleraumstände wegen« nur der
gleichzeitig abgegangene Brief an \textsc{Gauss} in dessen Hände, an dessen Schluss
\textsc{Wolfgang}, wie er an \textsc{Johann} schrieb, »eine kleine, klare Idee der Arbeit gab,
damit er nicht im voraus sich grause vor der Materie«. Der Sonderabzug
selbst kam nach längerer Zeit an \textsc{Wolfgang} zurück und ist Anfang Februar
1832 durch einen Bekannten der \textsc{Bolyais}, den in Göttingen studierenden
Baron v. \textsc{Zeyk}, \textsc{Gauss} übergeben worden (\textsc{Bol.} S. 91—92).
\par
Unter dem ersten Eindruck, den die Schrift auf ihn machte, schrieb \textsc{Gauss}
am 14. Februar 1832 an \textsc{Gerling}: »Noch bemerke ich, dass ich dieser Tage
eine kleine Schrift aus \textls[100]{Ungarn} über die nichteuklidische Geometrie erhalten
habe, worin ich alle \textls[100]{meine eigenen Ideen und Resultate} wiederfinde,
mit grosser Eleganz entwickelt, obwohl in einer für jemand, dem die Sache
fremd ist, wegen der Konzentrierung etwas schwer zu folgenden Form. Der
Verfasser ist ein \textls[100]{sehr} junger österreichischer Offizier, Sohn eines Jugendfreundes
von mir, mit dem ich 1798 mich oft über die Sache unterhalten
hatte, wiewohl damals meine Ideen noch viel weiter von der Ausbildung und
Reife entfernt waren, die sie durch das eigene Nachdenken dieses jungen
Mannes erhalten haben. Ich halte diesen jungen Geometer v. \textsc{Bolyai} für ein
Genie erster Grösse« (W. VIII, S. 220).
\par
Am 6. März folgte der wiederholt angeführte Brief an \textsc{Wolfgang}, in dem
\textsc{Gauss} seine Überraschung über das Zusammentreffen mit \textsc{Johann} ausdrückt
und bittet, diesen herzlich von ihm zu grüssen und ihm seine besondere Hochachtung
zu versichern (W. VIII, S 220—224). »\textsc{Gaussens} Antwort hinsichtlich
Deines Werkes«, schrieb \textsc{Wolfgang} an den Sohn, »ist sehr schön und gereicht
unserem Vaterlande und unserer Nation zur Ehre. Ein guter Freund sagt,
\clearpage\SourceLeaf{368}\ArithmeticHead{37}{R}{DIE GRUNDLAGEN DER GEOMETRIE.}\setcounter{footnote}{0}
\noindent es wäre eine grosse Satisfaktion« (\textsc{Bol.} S. 72). \textsc{Johann} selbst hat es als eine
grosse Enttäuschung und Kränkung empfunden, dass \textsc{Gauss} den \textit{Appendix}
keiner öffentlichen Anerkennung würdigte und das Vorrecht der ersten Entdeckung
für sich in Anspruch nahm (\textsc{Bol.} S. 95—97).
\par
Wie schon erwähnt wurde, gibt \textsc{Gauss} in dem Briefe als Probe ihm
eigentümlicher Untersuchungen einen einfachen Beweis für den Satz, dass in
der nichteuklidischen Geometrie der Inhalt des Dreiecks der Abweichung der
Winkelsumme von zwei Rechten proportional ist; der Umstand, dass damals
die Erinnerungen an den Verkehr mit \textsc{Wolfgang} in ihm wiederauftauchten,
macht es wahrscheinlich, dass er dabei an seine Untersuchungen vom September
1799 angeknüpft hatte. Er schliesst daran die Aufforderung, \textsc{Johann}
möge sich mit der entsprechenden Aufgabe für den Raum beschäftigen, nämlich
»den Kubikinhalt des Tetraeders (von vier Ebenen begrenzten Raumes)
zu bestimmen«. \textsc{Johann} hatte, wie sein Vater am 20. April 1835 an \textsc{Gauss}
schreibt (Br. G.-\textsc{Bolyai}, S. 115), die Auflösung der Aufgabe bereits ein Jahr
vor der Herausgabe des \textit{Appendix} gefunden. Der Nachlass \textsc{Johanns} enthält
in der Tat sogar mehrere Verfahren, die zur Lösung dienen können (\textsc{Bol.}
S. 109—118), darunter auch genau die Methode, die \textsc{Gauss} im Auge hatte
und die er, seiner oben (S. 17, 18) erwähnten Gewohnheit gemäss, bei der
Absendung des Briefes an \textsc{Wolfgang} vom 6. März 1832 in einem seiner Handbücher
angedeutet hat (W. VIII, S. 228).
\par
Auf das Volumen des Tetraeders bezieht sich noch eine zweite Aufzeichnung
von \textsc{Gauss}, die etwa aus dem Jahre 1841 stammt. Sie steht auf
einem Zettel, der sich in dem Sonderabdruck der Abhandlung \textsc{Lobatschefskijs}
vom Jahre 1836 über die Anwendung der imaginären Geometrie auf einige
Integrale gefunden hat; unter imaginärer Geometrie versteht der russische
Mathematiker die nichteuklidische Geometrie.
\par
\textsc{Lobatschefskij} (1793—1856) hatte in den Vorlesungen über Geometrie,
die er 1815/16 an der Universität Kasan hielt, noch ganz auf dem Boden der
euklidischen Geometrie gestanden und darin verschiedene Versuche zum Beweise
des Parallelenaxioms gemacht (\textsc{Lob.} S. 262, 378). Verraten schon diese
Vorlesungen eine eingehende Beschäftigung mit \textsc{Legendres} \textit{Elementen der Geometrie}
(\textsc{Lob.} S. 454), so lassen die späteren Schriften \textsc{Lobatschefskijs} erkennen,
dass er sich in den folgenden Jahren in tief eindringender Kritik mit \textsc{Legendre}
\clearpage\SourceLeaf{369}\ArithmeticHead{38}{L}{STÄCKEL, GAUSS ALS GEOMETER.}\setcounter{footnote}{0}
\noindent auseinandergesetzt und, indem er es wagte, Folgerungen aus der Annahme
des Nichtbestehens des Parallelenaxioms zu ziehen, sich allmählich mit dem
Gedanken von dessen Unbeweisbarkeit vertraut gemacht hat. Diesen Standpunkt
vertritt er in einem ungedruckt gebliebenen Lehrbuch der Geometrie vom
Jahre 1823 (\textsc{Lob.} S. 369). In den folgenden Jahren gelangte er zu der Erkenntnis,
dass es eine in sich widerspruchsfreie Geometrie gibt, die des Parallelenaxioms
nicht bedarf. Er entwickelte diese Geometrie soweit, dass er alle
ihre Aufgaben rein analytisch behandeln konnte; auch gab er allgemeine
Regeln zur Berechnung der Bogenlängen, Flächenräume und Rauminhalte.
Die Ergebnisse dieser Untersuchungen wurden am 12. Februar 1826 der Kasaner
Gelehrten Gesellschaft vorgelegt; veröffentlicht sind sie jedoch erst 1829
und 1830 im Kasaner Boten (\textsc{Lob.} S. 371). Ihnen folgte eine Reihe weiterer,
in russischer Sprache geschriebener Abhandlungen (1835—1838).
\par
Um seinen Gedanken Verbreitung im westlichen Europa zu verschaffen,
hatte \textsc{Lobatschefskij} 1837 in \textsc{Crelles} Journal eine kurze Darstellung seiner
imaginären Geometrie gegeben, die freilich zur Einführung in den Gegenstand
wenig geeignet war. \textsc{Gauss} scheint sie nicht beachtet zu haben, er ist vielmehr
wohl erst im Jahre 1840 auf \textsc{Lobatschefskij} aufmerksam geworden, als
dessen vortrefflich geschriebene deutsche Schrift: \textit{Geometrische Untersuchungen
zur Theorie der Parallellinien}, in \textsc{Gersdorfs} Repertorium abfällig besprochen
wurde (Brief an \textsc{Encke} vom 1. Febr. 1841, W. VIII, S. 232). Durch einen
merkwürdigen Zufall erhielt er um dieselbe Zeit durch den mit \textsc{Lobatschefskij}
befreundeten Physiker der Kasaner Universität \textsc{Knorr}, der ihn 1840 in Göttingen
besucht hatte, die schon erwähnte Abhandlung vom Jahre 1836. Später
hat ihm der Astronom W. \textsc{Struve} in Pulkowa die anderen, in den Kasaner
Gelehrten Schriften erschienenen Abhandlungen verschafft (W. VIII, S. 239);
woher \textsc{Gauss} die Abhandlung im Kasaner Boten vom Jahre 1829/30 bekommen
hat, ist unaufgeklärt (\textsc{Lob.} S. 435).
\par
Ein weiterer glücklicher Umstand war es, dass \textsc{Gauss} die in russischer
Sprache geschriebenen Schriften lesen konnte. »Die Aneignung irgend einer
neuen Fertigkeit als eine Art Verjüngung betrachtend« (Brief an \textsc{Schumacher}
vom 17. August 1839, Br. G.-\textsc{Sch.} III, S. 242) hatte er, nachdem er dem
Sanskrit keinen Geschmack abgewinnen konnte, im Frühjahr 1839 angefangen,
die russische Sprache zu erlernen. »Es dauerte kaum zwei Jahre, dass er ohne
\clearpage\SourceLeaf{370}\ArithmeticHead{39}{R}{DIE GRUNDLAGEN DER GEOMETRIE.}\setcounter{footnote}{0}
\noindent alle fremde Hilfe dieselbe so vollständig in seine Gewalt bekam, dass er nicht
nur alle Bücher in Prosa und Poesie mit Geläufigkeit lesen konnte, sondern
dass er sogar seine Korrespondenzen nach St. Petersburg mitunter in russischer
Sprache besorgte« (\textsc{Sartorius}, S. 91).
\par
Über die russischen Abhandlungen urteilt \textsc{Gauss} in dem Brief an \textsc{Gerling}
vom 8. Februar 1844, dass sie »mehr einem verworrenen Walde gleichen,
durch den es, ohne alle Bäume erst einzeln kennen gelernt zu haben, schwer
ist, einen Durchgang und Übersicht zu finden« (W. VIII, S. 237). Dagegen
lobt er die Konzinnität und Präzision der \textit{Geometrischen Untersuchungen} und
wiederholt dieses Lob in dem Brief an \textsc{Schumacher} vom 28. November 1846:
»Materiell für mich Neues habe ich .\ .\ .\ .\ nicht gefunden, aber die Entwicklung
ist auf anderm Wege gemacht, als ich selbst eingeschlagen habe,
und zwar von \textsc{Lobatschefskij} auf eine meisterhafte Art in echt geometrischem
Geiste. Ich glaube Sie auf das Buch aufmerksam machen zu müssen, welches
Ihnen gewiss ganz exquisiten Genuss gewähren wird« (W. VIII, S. 238).
\par
Auf einem Zettel, der sich in einem der beiden \textsc{Gauss} gehörenden Abdrücke
der \textit{Geometrischen Untersuchungen} vorgefunden hat, ist in gedrängter
Darstellung die bereits erwähnte (S. 31) Herleitung der Formeln der nichteuklidischen
Trigonometrie enthalten (W. VIII, S. 255—257); vermutlich ist
sie verfasst worden, als \textsc{Gauss} im Jahre 1846 »Veranlassung hatte, das Werkchen
wieder durchzusehen« (W. VIII, S. 238). Wenn dort (S. 31) bemerkt
wurde, dass die Aufzeichnung wohl den Gedankengang wiedergebe, den \textsc{Gauss}
im Jahre 1816 eingeschlagen hat, so muss hier hervorgehoben werden, dass
darin auch eine Auffassung zu Tage tritt, die \textsc{Gauss} erst später gewonnen hat.
Als Endergebnis werden nämlich Formeln erhalten, die mit den Gleichungen
der sphärischen Trigonometrie, bezogen auf eine Kugel vom Halbmesser \(1/k\),
identisch sind; die entsprechenden Gleichungen der nichteuklidischen Trigonometrie
folgen daraus, wenn der Konstanten \(k\) ein rein imaginärer Wert erteilt
wird. Dass diese Beziehung stattfindet, hatte \textsc{Lobatschefskij} am Schluss der
geometrischen Untersuchungen (S. 60) angemerkt. Sie erscheint bei ihm als
ein sonderbarer Zufall. Hat \textsc{Gauss} tiefer geschaut? Hat er durch den Buchstaben
\(k\) andeuten wollen, dass die beiden Geometrien dem allgemeineren
Begriff der Geometrie einer Mannigfaltigkeit konstanten Krümmungsmasses
untergeordnet werden können? Wie im fünften Abschnitt dieses Aufsatzes
\clearpage\SourceLeaf{371}\ArithmeticHead{40}{L}{STÄCKEL, GAUSS ALS GEOMETER.}\setcounter{footnote}{0}
\noindent dargelegt werden wird, spricht vieles dafür, die Frage zu bejahen. Dann aber
würde ein Licht fallen auf eine dunkle Stelle in dem vorher angeführten Brief
an \textsc{Schumacher} vom 28. November 1846: »Sie wissen, dass ich schon seit 54
Jahren (seit 1792) dieselbe Überzeugung habe (mit einer gewissen spätern Erweiterung,
deren ich hier nicht erwähnen will)« (W. VIII, S. 238). Darf man
noch weiter gehen? Hat \textsc{Gauss} seine ursprüngliche Überzeugung später dahin
erweitert, dass er den Geometrien, die sich je nach dem Vorzeichen des
Krümmungsmasses ergeben, volle Gleichberechtigung zubilligte, hat er dem
Gedanken \textsc{Riemanns} vorausgenommen, man brauche den Raum nur als unbegrenzte,
nicht als unendliche Mannigfaltigkeit aufzufassen? Die vorliegenden
Anhaltspunkte gestatten es nur, Vermutungen auszusprechen.
\par
Als \textsc{Johann Bolyai} am 3. November 1823 dem Vater von seinen neuen
Entdeckungen berichtet hatte (\textsc{Bol.} S. 85), ermahnte ihn dieser, sich mit der
Bekanntmachung zu beeilen, weil »manche Dinge gleichsam eine Epoche haben,
wo sie dann an mehreren Orten aufgefunden werden, gleichwie im Frühjahr
die Veilchen mehrwärts ans Licht kommen« (\textsc{Bol.} S. 86). Die Namen \textsc{Gauss},
\textsc{Schweikart}, \textsc{Taurinus}, \textsc{Lobatschefskij} sind ein Beweis dafür, wie richtig \textsc{Wolfgang}
geurteilt hatte.
\par
Man hat allerdings dieses Zusammentreffen dadurch seiner Merkwürdigkeit
zu entkleiden versucht, dass man vermutete, \textsc{Bolyai} und \textsc{Lobatschefskij}
verdankten \textsc{Gauss}, der ohne Zweifel als Erster sich von den Fesseln der Überlieferung
frei gemacht hat, die Fragestellung ihrer Untersuchungen (\textsc{Lob.} S. 428,
442). Dass \textsc{Schweikart} von \textsc{Gauss} unabhängig gewesen ist, unterliegt keinem
Zweifel; dagegen sind bei \textsc{Taurinus} Anregungen durch \textsc{Schweikart} und \textsc{Gauss}
wirksam gewesen, ohne dass ihm damit die Selbständigkeit in der Entdeckung
der nichteuklidischen Trigonometrie abgesprochen werden darf.
\par
Nachdem die hinterlassenen Schriften von \textsc{Gauss} und den beiden \textsc{Bolyai}
zugänglich geworden sind, können die Beziehungen zwischen ihnen als völlig
geklärt gelten; man beachte vor allem die beiden Tatsachen, dass \textsc{Wolfgang},
als \textsc{Gauss} im Herbst 1798 Göttingen verlassen hatte, das Parallelenaxiom zu
beweisen bemüht war, und dass, wie die S. 48 wiedergegebene Stelle aus einem
Briefe \textsc{Wolfgangs} beweist, \textsc{Johann} erst, nachdem er seine Untersuchungen
bereits begonnen hatte, von seinem Vater die Mitteilung erhielt, \textsc{Gauss} habe
\clearpage\SourceLeaf{372}\ArithmeticHead{41}{R}{DIE GRUNDLAGEN DER GEOMETRIE.}\setcounter{footnote}{0}
\noindent fruchtlos über die Parallelen nachgedacht, eine Warnung eher, denn eine
Anregung.
\par
Bei \textsc{Lobatschefskij} hat man an eine Vermittlung durch \textsc{Bartels} (1769—
1836) gedacht, der 1807 bis 1821 an der Universität Kasan gelehrt hat.
\textsc{Bartels} ist nämlich Hilfslehrer an der Schule gewesen, an der \textsc{Gauss} seinen
ersten Unterricht empfing und hat sich des Knaben hilfreich angenommen;
später, 1805 bis 1807, wo er als ein Schützling des Herzogs, wie \textsc{Gauss}, in
Braunschweig lebte, hat er mit diesem freundschaftlich verkehrt. Wenn aber
schon die ganze Entwicklung der Gedanken, wie sie vorher dargestellt worden
ist, für die volle Selbständigkeit \textsc{Lobatschefskijs} spricht, so kommt noch dazu,
dass \textsc{Bartels}, nach dem Zeugnis seines Schwiegersohnes O. \textsc{Struve}, in der
imaginären Geometrie mehr eine geistreiche Spekulation als ein die Wissenschaft
förderndes Werk gesehen hat; auch erinnert sich \textsc{Struve} nicht, dass
\textsc{Bartels} jemals von anklingenden Ideen bei \textsc{Gauss} gesprochen habe (\textsc{Lob.}
S. 378—382).
\clearpage\endgroup
